\documentclass[12pt]{article}
\usepackage[utf8]{inputenc}
\usepackage[T1]{fontenc}
\usepackage{hyperref}
\title{Kibiri Language Resources for K-12 Curriculum Development}
\author{Compiled by Claude Sonnet 5.5 (claude-sonnet-5-5) for Andr\'as Kornai}
\date{}
\begin{document}
\maketitle

\section{Introduction}
Kibiri (also Porome, Gibiri, Polome; ISO 639-3 \texttt{prm}) is a Papuan language of the Kikori District, Gulf Province, Papua New Guinea, with roughly 1,200 speakers (2011). It should not be confused with Rumu, which is also sometimes called Kibiri. Its classification is uncertain: it was treated as an isolate by Wurm, and Ross proposed a link to Kiwaian on thin evidence. My searches found only general descriptive pages and no teaching materials.

\section{Primary Educational Resources}
No primers, textbooks, or school programmes in Kibiri turned up in my searches. Nothing indicates that the language is taught in schools. A curriculum would have to be built from scratch together with the community.

\section{Linguistic and Reference Materials}
\begin{itemize}
\item Wikipedia, Porome language (overview, classification, speaker numbers): \url{https://en.wikipedia.org/wiki/Porome_language}
\item Endangered Languages Project record: \url{https://endangeredlanguages.com/node/98693}
\end{itemize}
I found no grammar, dictionary, or orthography document.

\section{Religious and Community Resources}
\begin{itemize}
\item Joshua Project profile of the Kibiri/Porome people group (mentions Tok Pisin as a language of wider communication): \url{https://joshuaproject.net/people_groups/12676/PP}
\end{itemize}
I found no scripture translation or community organisation page.

\section{Digital and Audio Resources}
\begin{itemize}
\item Wikitongues language page for \texttt{prm}: \url{https://wikitongues.org/languages/prm/}
\item Verbix endangered-languages page: \url{https://docs.verbix.com/EndangeredLanguages/KibiriPorome}
\end{itemize}
I did not verify that either page contains audio.

\section{Curriculum Development Context}
Very little published material exists, and the sources I found disagree on vitality: Wikitongues gives EGIDS 6a (vigorous) while another source calls the language threatened. Any curriculum would need fieldwork and community-led documentation, and would probably begin with oral and literacy-bridging materials alongside Tok Pisin and English.

\end{document}